% !TEX encoding = UTF-8 Unicode
\section{Qualitätsanforderungen}


\newcommand{\haupt}[1]{\textbf{#1}		&				& 			&			& 	 			    \\}
\newcommand{\noteA}[1]{\hline#1			&		X		& 			&			& 	 			    \\}
\newcommand{\noteB}[1]{\hline#1			&		 		& 	   X		&			& 	 			    \\}
\newcommand{\noteC}[1]{\hline#1			&				& 			&	   X		& 	 			    \\}
\newcommand{\noteD}[1]{\hline#1			&				& 			&			& 	 	X		    \\}
\newcommand{\noteX}[1]{\hline#1			&		-		& 	    -		&	    -     	& 	 	-		    \\}

\neupunkt{QF}
Die Qualitätsanforderungen sind in der Tabelle \ref{quali} aufgeführt.
\punkt{QF} Alle Reaktionszeiten auf Benutzeraktionen müssen unter 100ms liegen. Die maximale Verarbeitungslatenz des Systems, sollte die Maximums-Latenz aller vom System enthaltenden Plugins, nicht überschreiten.

\begin{table}[htdp]
\caption{Qualitätsanforderungen an VSTForx}
\begin{center}
\begin{tabular}{c|c|c|c|c}

Systemqualität				&	sehr gut	& 	gut		&	normal	& 	nicht relevant \\
\hline 
\hline 
\haupt{Funktionalität}	
\noteA{Angemessenheit} 	
\noteC{Genauigkeit}				
\noteA{Interoperabilität}
\noteD{Sicherheit}
\noteB{Konformität}

\hline
\hline 
\haupt{Zuverlässigkeit}
\noteA{Reife}
\noteB{Fehlertoleranz}
\noteD{Wiederherstellbarkeit}
\noteB{Konformität}

\hline 
\hline
\haupt{Benutzbarkeit}
\noteC{Verständlichkeit}
\noteB{Erlernbarkeit}
\noteA{Bedienbarkeit}
\noteC{Attraktivität}
\noteB{Konformität}

\hline 
\hline
\haupt{Effizienz}
\noteA{Zeitverhalten}
\noteA{Verbrauchsverhalten}
\noteA{Konformität}

\hline 
\hline
\haupt{Wartbarkeit}
\noteC{Analysierbarkeit}
\noteB{Änderbarkeit}
\noteA{Stabilität}
\noteA{Testbarkeit}
\noteB{Konformität}

\hline 
\hline
\haupt{Portabilität}
\noteA{Anpassbarkeit}
\noteA{InstalIierbarkeit}
\noteA{Koexistenz}
\noteD{Austauschbarkeit}
\noteA{Konformität}



\end{tabular}
\end{center}
\label{quali}
\end{table}
